\documentclass[12pt,a4paper]{ctexart}

% === 页面设置 ===
\usepackage[top=2.5cm,bottom=2.5cm,left=2.5cm,right=2.5cm]{geometry}

% === 数学和物理 ===
\usepackage{amsmath,amssymb,amsthm}
\usepackage{physics}
\usepackage{slashed}

% === 图形和表格 ===
\usepackage{graphicx}
\usepackage{float}
\usepackage{booktabs}
\usepackage{array}
\usepackage{caption}

% === 引用 ===
\usepackage{hyperref}
\usepackage[numbers,sort&compress]{natbib}

% === 其他 ===
\usepackage{xcolor}
\usepackage{enumitem}
\usepackage{fancyhdr}
\usepackage{multirow}
\usepackage{tikz}

\hypersetup{
    colorlinks=true,
    linkcolor=blue,
    citecolor=blue,
    urlcolor=blue,
}

\theoremstyle{definition}
\newtheorem{definition}{定义}[section]
\newtheorem{remark}{注记}[section]
\newtheorem{theorem}{定理}[section]

\title{\heiti 格点QCD中的光锥与光前：\\
从Minkowski光锥关联到Euclidean格点计算}
\author{基于本库文献的综合解析}
\date{\today}

\begin{document}

\maketitle

\begin{abstract}
\noindent
部分子物理的核心可观测量——部分子分布函数（PDFs）、广义部分子分布（GPDs）、分布振幅（DAs）和横向动量依赖分布（TMDs）——本质上是由Minkowski时空中的光锥（light-cone）关联函数定义的。然而，格点QCD作为从第一性原理求解QCD非微扰性质的主要工具，是在Euclidean时空中表述的，无法直接处理实时间依赖的光锥关联。这一基本矛盾长期阻碍了部分子物理的格点计算。大动量有效理论（LaMET）和赝PDF方法的提出彻底改变了这一局面：通过Lorentz boost，将空间方向的等时关联函数在大动量极限下映射到光锥关联函数，从而在Euclidean格点上间接计算光锥物理量。本文系统梳理了这一理论框架的完整逻辑链条：从光锥坐标和光前量子化的基本概念出发，经过Feynman的无穷大动量系图像，到LaMET的有效场论表述和因子化定理的严格证明，再到准PDF、赝PDF、匹配核、混合重整化等关键技术的发展，最终展示这一框架在胶子PDF、螺旋度分布、介子光锥振幅等实际计算中的应用。
\end{abstract}

\tableofcontents
\newpage

% ============================================================
\section{光锥与光前：基本概念}
% ============================================================

\subsection{光锥坐标}

在狭义相对论中，光锥（light-cone）坐标是一组特别适合描述以接近光速运动的物理系统的坐标。对于沿$z$方向运动的物理系统，定义\cite{Ji2014LaMET}：

\begin{equation}
\xi^\pm = \frac{1}{\sqrt{2}}(\xi^0 \pm \xi^3),
\label{eq:lightcone_coords}
\end{equation}

其中$\xi^0 = t$是物理时间，$\xi^3 = z$是空间坐标。$\xi^+$被称为``光前时间''（light-front time），$\xi^-$被称为``光前纵向坐标''。横向坐标$\vec{\xi}_\perp = (\xi^1, \xi^2)$不变。相应的动量分量为$P^\pm = (P^0 \pm P^3)/\sqrt{2}$。

光锥坐标的关键性质包括\cite{Ji2014LaMET}：

\begin{enumerate}
    \item \textbf{度规：}在$(\xi^+, \xi^-, \xi^1, \xi^2)$坐标下，$g_{+-} = g_{-+} = 1$，$g_{11} = g_{22} = -1$，所有其他分量为零。
    \item \textbf{标量积：}$a \cdot b = a^+ b^- + a^- b^+ - \vec{a}_\perp \cdot \vec{b}_\perp$。
    \item \textbf{Lorentz boost：}沿$z$方向的boost简化为$\xi^+ \to e^{\omega}\xi^+$，$\xi^- \to e^{-\omega}\xi^-$，其中$\omega$是快度参数。
\end{enumerate}

\subsection{光前量子化与无穷大动量系}

光前量子化的思想可以追溯到Dirac在1949年的开创性工作\cite{Ji2014LaMET}。其基本思路是：将$\xi^+$视为新的``时间''坐标，引入等``时间''对易关系来量子化场论。在$\xi^+ = 0$处的场具有按Fock粒子展开的平面波形式。新的哈密顿量为$\hat{H}_{\text{LC}} = \hat{P}^-$，可以用于发展``时间无关''的微扰理论。

在光前量子化框架下，最方便的是选择$A^+ = 0$规范（光锥规范）。此时Fock粒子就是部分子（parton）。强子态可以用Fock展开来表示\cite{Ji2014LaMET}：

\begin{equation}
|P\rangle = \sum_{n\alpha,\lambda_i} \int \prod_i \frac{dx_i d^2k_{\perp,i}}{\sqrt{2x_i}(2\pi)^3} \psi_{n\alpha}(x_i, k_{\perp,i}, \lambda_i) |n\alpha: x_i P^+, k_{\perp,i}, \lambda_i\rangle,
\label{eq:Fock_expansion}
\end{equation}

其中$n$是部分子数目，$\lambda_i$是螺旋度标记，$x_i$是纵向动量份额，$k_{\perp,i}$是横向动量。

\textbf{无穷大动量系（Infinite Momentum Frame, IMF）}是Feynman在1969年提出的等价图像\cite{Ji2014LaMET}：考虑一个以接近光速运动的强子，其动量$P^z \to \infty$。在这个极限下，强子可以被视为一束互不相互作用的部分子，它们由动量密度$q(x)$（夸克）和$g(x)$（胶子）来表征，其中$x = k^z/P^z$是部分子携带的纵向动量份额。

IMF与光锥形式之间的联系在于：IMF中的等时关联函数在无穷大boost下等价于光锥关联函数\cite{Ji2013Euclidean, Ji2014LaMET}。如季向东所指出的\cite{Ji2013Euclidean}：

\begin{quote}
``Intuitively, the parton distribution can be regarded for the nucleon in the rest frame but the probe is traveling at the speed of light, and hence the light-cone correlation; On the other hand, the distribution is obtained from a static probe on a nucleon traveling at the speed of light.''
\end{quote}

\section{光锥关联函数与部分子物理}

\subsection{PDF的光锥定义}

在现代QCD因子化定理中\cite{Izubuchi2018}，部分子分布函数由光锥关联算符的核子矩阵元来定义。例如，裸的（bare）非极化夸克PDF为\cite{Izubuchi2018}：

\begin{equation}
q(x, \epsilon) \equiv \int \frac{d\xi^-}{4\pi} e^{-ixP^+\xi^-} \langle P| \bar{\psi}(\xi^-) \gamma^+ U(\xi^-, 0) \psi(0) |P\rangle,
\label{eq:lightcone_quark_pdf}
\end{equation}

其中核子动量$P^\mu = (P^0, 0, 0, P^z)$，Wilson线

\begin{equation}
U(\xi^-, 0) = \mathcal{P} \exp\left[-ig \int_0^{\xi^-} d\eta^- A^+(\eta^-)\right]
\label{eq:wilson_line}
\end{equation}

保证了非局域双线性算符的规范不变性。在$A^+ = 0$规范下，$U = 1$，PDF简化为简单的Fourier变换。

\textbf{胶子PDF}的光锥定义在伴随表示下为\cite{Chen2025gluon}：

\begin{equation}
xg(x, \mu) = \int \frac{dz^-}{2\pi P^+} e^{-ixP^+z^-} \langle P| F^a_{+\mu}(z^-) U^{ab}(z^-, 0) F^{b\mu}_{\phantom{b\mu}+}(0)(\mu) |P\rangle,
\label{eq:lightcone_gluon_pdf}
\end{equation}

其中$F^{\mu\nu}$是规范场张量，$U^{ab}$是伴随表示中的Wilson线。

\textbf{极化胶子分布}（胶子螺旋度）则由对偶场强张量$\tilde{F}^{\mu\nu} = \frac{1}{2}\epsilon^{\mu\nu\rho\sigma}F_{\rho\sigma}$定义\cite{Egerer2022}：

\begin{equation}
\Delta G = \int dx\, \frac{i}{2xP^+} \int \frac{d\xi^-}{2\pi} e^{-ixP^+\xi^-} \langle PS| F^{+\alpha}_a(\xi^-) W^{ab}(\xi^-, 0) \tilde{F}^+_{\alpha,b}(0) |PS\rangle。
\label{eq:DeltaG_lightcone}
\end{equation}

\subsection{为什么格点QCD难以直接计算光锥关联？}

Wilson的格点QCD方法\cite{Ji2013Euclidean}是在离散化的Euclidean时空中表述的，其中时间轴被Wick旋转为虚时间：$t \to -i\tau$。这带来了一个根本性的困难：

\begin{quote}
``Since the lattice theory is defined in a discretized Euclidean space with imaginary time, it is very difficult to calculate Minkowskian quantities with real-time dependence such as the PDF.''\cite{Izubuchi2018}
\end{quote}

光锥关联函数涉及沿$\xi^- = (t-z)/\sqrt{2}$方向分离的场，不可避免地包含实时间依赖性。在Euclidean时空中，人们只能计算等时（$\tau = 0$）关联函数或局域算符的矩阵元。

传统的解决方法是计算PDF的\textbf{矩}（moments）\cite{Ji2013Euclidean}：

\begin{equation}
\int dx\, x^{n-1} q(x, \mu^2) = a_n(\mu^2),
\label{eq:moments}
\end{equation}

其中$a_n$是局域twist-2算符的约化矩阵元：

\begin{equation}
\langle P| \mathcal{O}^{\mu_1\ldots\mu_n}(\mu^2) |P\rangle = 2a_n(\mu^2)(P^{\mu_1}\ldots P^{\mu_n} - \text{trace})。
\label{eq:local_operators}
\end{equation}

局域算符为$\mathcal{O}^{\mu_1\ldots\mu_n} = \bar{\psi}\gamma^{(\mu_1}iD^{\mu_2}\ldots iD^{\mu_n)}\psi - \text{trace}$。然而，由于格点离散化破坏了$O(4)$旋转对称性，高维算符之间存在复杂的混合，使得实际上只能计算前两三个矩\cite{Ji2013Euclidean}。

\section{LaMET：大动量有效理论}

\subsection{核心思想：用Lorentz Boost连接Euclidean与光锥}

季向东在2013年提出了一个革命性的方案\cite{Ji2013Euclidean}：利用Lorentz boost，将空间方向的等时关联函数在大动量极限下映射到光锥关联函数。

核心论证如下\cite{Ji2013Euclidean}：考虑一个空间分离的非局域夸克算符

\begin{equation}
\mathcal{O}^z = \bar{\psi}\gamma^z iD^z\ldots iD^z\psi - \text{trace},
\label{eq:spatial_operator}
\end{equation}

其中所有Lorentz指标都取为$z$方向。在大动量$P^z$的核子态中计算其矩阵元。根据Lorentz对称性，其矩阵元为：

\begin{equation}
\langle P| \bar{\psi}\gamma^z iD^z\ldots iD^z\psi |P\rangle = 2a_n(\mu^2)(P^z)^n \times \left[1 + \mathcal{O}\left(\frac{\Lambda_{\text{QCD}}^2}{(P^z)^2}, \frac{M^2}{(P^z)^2}\right)\right]。
\label{eq:spatial_ME}
\end{equation}

关键的结论是：在大$P^z$极限下，非领头项被幂次压低。将上述结果逆推为非局域形式，得到\textbf{准PDF}（quasi-PDF）\cite{Ji2013Euclidean}：

\begin{definition}[准夸克分布]
\begin{equation}
q(x, \mu^2, P^z) = \int \frac{dz}{4\pi} e^{izk^z} \langle P| \bar{\psi}(z) \gamma^z \exp\left[-ig\int_0^z dz' A^z(z')\right] \psi(0) |P\rangle + \mathcal{O}\left(\frac{\Lambda^2}{(P^z)^2}, \frac{M^2}{(P^z)^2}\right),
\label{eq:quasi_pdf}
\end{equation}
其中$x = k^z/P^z$。
\end{definition}

\textbf{几何直觉}\cite{Ji2013Euclidean, Ji2014LaMET}：考虑连接$(0,0,0,z)$与坐标原点的空间线段。当boost速度趋近光速时，这条类空线段被倾斜到光锥方向。它不可能严格位于光锥上（因为boost不改变不变长度），但这种微小的偏离光锥性只引入在$P^z \to \infty$极限下消失的幂次修正。

\begin{figure}[H]
\centering
\begin{tikzpicture}[scale=1.5]
    % Axes
    \draw[->] (-0.5,0) -- (6,0) node[right] {$\xi^3 = z$};
    \draw[->] (0,-0.5) -- (0,5) node[above] {$\xi^0 = t$};

    % Light cone
    \draw[dashed, gray] (0,0) -- (5,5) node[right] {$\xi^+ = 0$};
    \draw[dashed, gray] (0,0) -- (5,-5) node[right] {};

    % Space-like segment at rest (horizontal)
    \draw[thick, blue, ->] (0,0) -- (2,0) node[below] {空间分离$z$（静止系）};

    % Boosted segment
    \draw[thick, red, ->] (0,0) -- (1,2.5) node[above left] {大$P^z$下的boost};

    % Light-cone direction
    \draw[thick, green!60!black, ->] (0,0) -- (1.77,1.77) node[right] {光锥方向};

    \node[font=\small] at (3, 3.5) {$\xi^+ = (\xi^0 + \xi^3)/\sqrt{2}$};
    \node[font=\small] at (3, -1.5) {$\xi^- = (\xi^0 - \xi^3)/\sqrt{2}$};
\end{tikzpicture}
\caption{空间分离线段在Lorentz boost下趋近光锥方向的示意图}
\label{fig:boost_to_lightcone}
\end{figure}

\subsection{LaMET的有效场论表述}

季向东随后将上述思想系统地表述为\textbf{大动量有效理论（LaMET）}的框架\cite{Ji2014LaMET}。LaMET与重夸克有效理论（HQET）之间存在精确的类比：

\begin{center}
\begin{tabular}{lcc}
\toprule
& HQET & LaMET \\
\midrule
大标度        & 重夸克质量 $m_b$ & 强子动量 $P^z$ \\
有效理论      & $m_b \to \infty$ & $P^z \to \infty$ \\
小参量展开    & $1/m_b$          & $1/P^z$ \\
匹配因子      & $Z(m_b/\Lambda, \Lambda/\mu)$ & $Z(P^z/\Lambda, \Lambda/\mu)$ \\
有效理论观测量 & $o(\mu)$         & $f(\mu)$（光锥量） \\
\bottomrule
\end{tabular}
\end{center}

LaMET的基本展开公式为\cite{Ji2014LaMET}：

\begin{equation}
F(P^z/\Lambda) = Z(P^z/\Lambda, \Lambda/\mu) f(\mu) + \mathcal{O}(1/(P^z)^2) + \ldots,
\label{eq:LaMET_expansion}
\end{equation}

其中$F$是格点上可计算的准观测量（依赖于大强子动量$P^z$），$f(\mu)$是光锥观测量（$P^z \to \infty$的有效理论），$Z$是微扰可计算的匹配因子。

\textbf{因子化定理}在2018年由Izubuchi等人给出了严格证明\cite{Izubuchi2018}。准PDF与光锥PDF之间的关系是：

\begin{equation}
\tilde{x}\tilde{q}(x, P^z) = \int_{-\infty}^{\infty} \frac{dy}{|y|} C\left(\frac{x}{y}, \frac{\mu}{yP^z}\right) q(y, \mu) + \mathcal{O}\left(\frac{\Lambda_{\text{QCD}}^2}{(xP^z)^2}, \frac{\Lambda_{\text{QCD}}^2}{((1-x)P^z)^2}\right),
\label{eq:factorization}
\end{equation}

其中$C$是微扰匹配核，在$\alpha_s$中具有微扰展开：

\begin{equation}
C(\xi, \mu/P^z) = \delta(1-\xi) + \frac{\alpha_s}{2\pi} C^{(1)}(\xi, \mu/P^z) + \ldots。
\label{eq:matching_kernel}
\end{equation}

\textbf{极限交换的非对易性}\cite{Ji2013Euclidean, Ji2014LaMET}：

\begin{quote}
``In a typical Feynman diagram, the order of regularization UV divergences vs.\ $P^z \to \infty$ does not commute with each other. While defining the parton distribution requires $P^z$ to be much larger than the cut-off scale, the lattice simulations only makes sense when the UV cut-off is much larger than the hadron momentum.''
\end{quote}

这一观察是理解LaMET因子化结构的关键：光锥PDF是在取$P^z \to \infty$之后再进行重整化定义的；而格点上的准PDF是在先施加格点正规化（UV截断）之后、在有限$P^z$下计算的。两者之间通过微扰匹配核$Z$来联系，$Z$中包含了所有与两种极限顺序差异相关的UV物理。

\subsection{准光前关联、准分布与赝分布}

在LaMET框架的发展过程中，出现了三种密切联系但又有所区别的计算方案\cite{Yao2023}：

\begin{enumerate}
    \item \textbf{准光前（quasi-LF）关联：}定义在Euclidean空间间隔上的等时夸克和胶子关联函数。通过\textbf{大动量因子化}（动量空间）或\textbf{短距离因子化}（坐标空间）与光前关联函数建立联系\cite{Yao2023}。

    \item \textbf{准分布（quasi-PDF）：}由准光前关联函数通过Fourier变换得到的动量空间分布：
    \begin{equation}
    \tilde{x}\tilde{q}(x, P^z) = \frac{1}{\mathcal{N}} \int \frac{dz}{2\pi P^z} e^{-ixzP^z} \tilde{h}^R(z, P^z),
    \label{eq:quasi_pdf_fourier}
    \end{equation}
    其中$\tilde{h}^R(z, P^z) = \langle P|O(z)|P\rangle_R$是重整化矩阵元。这是最初的LaMET方案\cite{Ji2013Euclidean}。

    \item \textbf{赝分布（pseudo-PDF）：}由Radyushkin提出的方案，通过引入Ioffe-时间分布（ITD）来避免Fourier变换中的大$z$困难\cite{Izubuchi2018}。Ioffe时间$\nu \equiv -(p \cdot z)$是DIS截面分析中原始变量的格点类比。赝ITD $\tilde{I}_p(\nu, z^2)$与光锥ITD $I_p(\nu, \mu^2)$通过坐标空间的因子化公式联系：
    \begin{equation}
    \tilde{I}_p(\nu, z^2) = \int_0^1 du\, C(u, \mu^2 z^2) I_p(u\nu, \mu^2) + \mathcal{O}(z^2\Lambda_{\text{QCD}}^2)。
    \label{eq:pseudo_factorization}
    \end{equation}
\end{enumerate}

Izubuchi等人\cite{Izubuchi2018}严格证明了\textbf{准PDF方法和赝PDF方法遵循等价的因子化公式}。两种方法的核心区别在于因子化是在动量空间还是在坐标空间中执行的。

\subsection{胶子准PDF的构造}

对于胶子，准PDF的构造比夸克情形更为复杂，因为胶子场强张量$F^{\mu\nu}$具有两个Lorentz指标\cite{Balitsky2020,NoteGluonPDFs,Chen2025gluon}。胶子算符的可乘性重整化形式为\cite{Chen2025gluon,Zhang2019}：

\begin{equation}
O(z) = M_{tx;tx}(z) + M_{ty;ty}(z) - 2M_{xy;xy}(z),
\label{eq:gluon_operator}
\end{equation}

其中在基础表示下\cite{Chen2025gluon}：

\begin{equation}
M_{\mu\lambda;\nu\rho}(z) = \sum_{\vec{x}} \operatorname{Tr}[F_{\mu\lambda}(\vec{x} + z\hat{z}) U(\vec{x} + z\hat{z}, \vec{x}) F_{\nu\rho}(\vec{x}) U(\vec{x}, \vec{x} + z\hat{z})]。
\label{eq:gluon_ME}
\end{equation}

这一组合的选择基于：$M_{ti;it}$和$M_{ij;ji}$具有相同的单圈紫外反常量纲，使得两者之差在至少单圈层次上是可乘性重整化的\cite{Balitsky2020}。

\textbf{Euclidean量与Minkowski量的转换}对于胶子算符至关重要\cite{NoteGluonPDFs}。在Euclidean时空中，场强张量的$tx$、$ty$分量满足：

\begin{equation}
F^{tx,(M)}F^{(M)}_{tx} = -F^{(E)}_{tx}F^{(E)}_{tx}, \qquad
F^{xy,(M)}F^{(M)}_{xy} = F^{(E)}_{xy}F^{(E)}_{xy}。
\label{eq:ME_conversion}
\end{equation}

这就是胶子流算符中为什么出现``$-$''号的原因\cite{NoteGluonPDFs}：

\begin{equation}
\begin{aligned}
O^{(0)}_g(z, t_i) &= \left[ F^{0i}(z,t_i)W[z,0]F_{0i}(0,t_i)W[0,z] + 2F^{12}(z,t_i)W[z,0]F_{12}(0,t_i)W[0,z] \right]_{\text{Mink}} \\
&= \left[ -F^{0i}(z,t_i)W[z,0]F_{0i}(0,t_i)W[0,z] + 2F^{12}(z,t_i)W[z,0]F_{12}(0,t_i)W[0,z] \right]_{\text{Eucl}}。
\end{aligned}
\label{eq:gluon_current}
\end{equation}

\section{重整化：从裸格点矩阵元到连续物理量}

\subsection{裸矩阵元中的发散结构}

裸准光前矩阵元$\tilde{h}^B(z, P^z, 1/a)$同时包含多种发散\cite{Ji2021hybrid,Chen2025gluon}：

\begin{itemize}
    \item \textbf{线性发散：}来自Wilson线的自能，形式为$e^{-\delta m |z|}$，其中$\delta m \sim 1/a$；
    \item \textbf{对数紫外发散：}来自算符本身的 anomalous dimension；
    \item \textbf{对数红外发散：}在非零动量下需要考虑。
\end{itemize}

\subsection{混合重整化方案}

为解决重整化问题，季向东等人提出了\textbf{混合重整化方案}\cite{Ji2021hybrid}。其核心思想是根据距离尺度采用不同的重整化策略：

\begin{enumerate}
    \item \textbf{短距离（$|z| < z_s$）：}通过除以零动量裸矩阵元来消除发散：
    \begin{equation}
    \tilde{h}^R_{\text{short}}(z, P^z) = \frac{\tilde{h}^n_B(z, P^z, 1/a)}{\tilde{h}^n_B(z, P^z=0, 1/a)}。
    \end{equation}

    \item \textbf{长距离（$|z| > z_s$）：}采用自重整化程序，重整化因子包含线性发散和反对数的重求和：
    \begin{equation}
    \begin{aligned}
    Z_R(z, 1/a, \mu) = \exp\Bigg\{&\frac{kz}{a}\ln[a\Lambda_{\text{QCD}}] + m_0 z \\
    &+ \frac{5C_A}{3b_0} \ln\left[\frac{\ln(1/a\Lambda_{\text{QCD}})}{\ln(\mu/\Lambda_{\text{QCD}})}\right] \\
    &+ \frac{1}{2}\ln\left[\left(1 + \frac{d}{\ln[a\Lambda_{\text{QCD}}]}\right)^2\right]\Bigg\}。
    \end{aligned}
    \label{eq:ZR}
    \end{equation}
\end{enumerate}

完整的重整化矩阵元为\cite{Chen2025gluon}：

\begin{equation}
\tilde{h}^R(z, P^z) = \frac{\tilde{h}^n_B(z, P^z, 1/a)}{\tilde{h}^n_B(z, P^z=0, 1/a)} \theta(z_s - |z|) + \eta_s \frac{\tilde{h}^n_B(z, P^z, 1/a)}{Z_R(z, 1/a)} \theta(|z| - z_s),
\label{eq:hybrid_renorm}
\end{equation}

其中$\eta_s$是方案转换因子，确保重整化关联函数在$z = z_s$处的连续性。

\subsection{Ioffe-时间分布与大$\lambda$外推}

重整化矩阵元$\tilde{h}^R(z, P^z)$是Ioffe-时间$\lambda = zP^z$的函数。为了进行Fourier变换得到动量空间的PDF，需要所有$\lambda$值处的矩阵元。然而，格点数据只能在有限的$z$和$P^z$范围内获取，且在大$\lambda$区域信噪比急剧恶化\cite{Chen2025gluon}。

解决方法是进行\textbf{大$\lambda$外推}\cite{Chen2025gluon}：

\begin{equation}
\tilde{h}^R(\lambda) = l_1 \lambda^{-a_1} e^{-\lambda/\lambda_0},
\label{eq:lambda_extrap}
\end{equation}

其中$\lambda^{-a_1}$与非极化PDF在端点区域的幂律行为相关，指数因子反映了有限动量下关联长度$\lambda_0$有限的物理预期。

\section{LaMET框架的进一步推广}

\subsection{广义部分子分布（GPDs）}

GPDs是PDFs的非前行推广，描述在初态和末态具有不同动量的核子之间的部分子关联。在LaMET框架中，GPD $E(x,\xi,t)$和$H(x,\xi,t)$可以从以下矩阵元计算\cite{Ji2013Euclidean}：

\begin{equation}
F(x, \xi, t, P^z) = \int \frac{dz}{2\pi} e^{-izk^z} \langle P'| \bar{\psi}(-z/2) \gamma^z L(-z/2, z/2) \psi(z/2) |P\rangle,
\label{eq:GPD_quasi}
\end{equation}

其中$\vec{P}' + \vec{P} = 2(0,0,P^z)$，在$P^z \to \infty$、$\xi = (P^{z\prime} - P^z)/P^z$固定、$x = k^z/P^z$固定的极限下。

Yao、Ji和Zhang\cite{Yao2023}系统推导了味单态和非单态组合在非前行运动学下的完整匹配核，为从格点QCD提取所有leading-twist GPDs提供了完整的理论手册。

\subsection{横向动量依赖分布（TMDs）}

TMDs描述部分子纵向动量和横向动量的联合分布。在LaMET中\cite{Ji2013Euclidean}：

\begin{equation}
q(x, k_\perp, P^z, \mu^2) = \int \frac{dz}{4\pi} d^2\vec{r}_\perp e^{i(zk^z + \vec{k}_\perp\cdot\vec{r}_\perp)} \langle P| \bar{\psi}(\vec{r}_\perp, z) L^\dagger(\pm\infty; (\vec{r}_\perp, z)) \gamma^z L(\pm\infty; 0) \psi(0) |P\rangle。
\label{eq:TMD_quasi}
\end{equation}

与准PDF不同，TMD的$P^z$依赖性包含大的对数$\ln P^z$，在$P^z \to \infty$极限下不消失\cite{Ji2013Euclidean}。因此，在TMD情形中需要推导关于$P^z$的演化方程（Collins-Soper方程）。

\subsection{分布振幅（DAs）}

$\pi^+$介子的leading光锥分布振幅可以从以下矩阵元计算\cite{Ji2013Euclidean}：

\begin{equation}
\phi_0(x, P^z) 2P^z = \int \frac{dz}{2\pi} e^{izk^z} \langle 0| \bar{d}(0) \gamma^z \gamma_5 L(0, z) u(z) |\pi^+(P)\rangle。
\label{eq:DA_quasi}
\end{equation}

Yao等人\cite{Yao2023}给出的统一框架包含了DAs作为非前行运动学的$P^{\prime} \to 0$极限情形。

\subsection{Wigner分布}

Wigner分布是GPDs和TMDs的生成函数。在LaMET中\cite{Ji2013Euclidean}：

\begin{equation}
\begin{aligned}
W(x, k_\perp, b_\perp, P^z, \mu^2) &= \int \frac{dz}{4\pi} d^2\Delta_\perp d^2\vec{r}_\perp e^{-i\Delta_\perp\cdot b_\perp} e^{i(zk^z + \vec{k}_\perp\cdot\vec{r}_\perp)} \\
&\quad \times \langle P'| \bar{\psi}(\vec{r}_\perp, z) L^\dagger(\pm\infty; (\vec{r}_\perp, z)) \gamma^z L(\pm\infty; 0) \psi(0) |P\rangle。
\end{aligned}
\label{eq:Wigner_quasi}
\end{equation}

\subsection{光锥波函数}

LaMET框架同样适用于光锥波函数的计算\cite{Ji2014LaMET}。以$\pi^+$介子为例：

\begin{equation}
\hat{F}_\pm(z, \xi_\perp) = \bar{d}(0) W^\dagger(\pm\infty, 0; 0) \gamma^z \gamma_5 W(\pm\infty, z; \xi_\perp) u(z, \xi_\perp),
\label{eq:LCWF_quasi}
\end{equation}

其中$W(\pm\infty, z; \xi_\perp) = \mathcal{P}\exp[\mp ig\int_z^{\pm\infty} dz' A^z(z', \xi_\perp)]$是空间方向的Wilson线。

\section{光锥物理的格点实践}

\subsection{连续极限和外推}

在实际格点计算中，准PDF同时依赖于$P^z$和格点间距$a$。要获得物理的光锥PDF，需要进行联合外推\cite{Chen2025gluon}：

\begin{equation}
xg(x, P^z, a) = xg_0(x) + a^2 f(x) + a^2 (P^z)^2 h(x) + \frac{d(x)}{(P^z)^2},
\label{eq:joint_extrap}
\end{equation}

其中各项分别代表：
\begin{itemize}
    \item $xg_0(x)$：外推至连续极限和无穷大动量极限的物理PDF；
    \item $a^2 f(x)$：与动量无关的离散化误差；
    \item $a^2 (P^z)^2 h(x)$：依赖于动量的离散化误差；
    \item $d(x)/(P^z)^2$：领头的高阶twist（幂次）修正。
\end{itemize}

\subsection{匹配核的具体形式}

对于非极化胶子PDF，混合方案下的NLO匹配核由ratio方案核加上混合修正项组成\cite{Chen2025gluon,Balitsky2020,Yao2023}：

\begin{equation}
C^{\text{ratio}}\left(\xi, \frac{\mu}{yP^z}\right) = \delta(1-\xi) + \frac{\alpha_s(\mu)C_A}{2\pi}
\begin{cases}
\left[\frac{2(1-\xi+\xi^2)^2}{1-\xi}\ln\frac{\xi}{\xi-1} + \frac{11-28\xi+18\xi^2-12\xi^3}{6(1-\xi)}\right]_+ & \xi > 1 \\
\left[\frac{2(1-\xi+\xi^2)^2}{1-\xi}\left(-\ln\frac{\mu^2}{4y^2(P^z)^2} + \ln(\xi(1-\xi))\right) - \cdots\right]_+ & 0 < \xi < 1 \\
\left[-\frac{2(1-\xi+\xi^2)^2}{1-\xi}\ln\frac{\xi}{\xi-1} - \cdots\right]_+ & \xi < 0。
\end{cases}
\label{eq:Cratio_full}
\end{equation}

混合修正项为\cite{Chen2025gluon}：

\begin{equation}
C^{\text{hybrid}}\left(\xi, \lambda_s, \frac{\mu}{yP^z}\right) = C^{\text{ratio}}\left(\xi, \frac{\mu}{yP^z}\right) + \frac{\alpha_s(\mu)C_A}{2\pi} \frac{5}{6}\left[-\frac{1}{|1-\xi|} + \frac{2\,\text{Si}((1-\xi)|y|\lambda_s)}{\pi(1-\xi)}\right]_+。
\label{eq:Chybrid}
\end{equation}

\subsection{从格点结果到物理预言的完整流程}

总结LaMET框架下从格点QCD计算到物理PDF的完整流程\cite{Chen2025gluon,Ji2021hybrid}：

\begin{enumerate}
    \item \textbf{计算裸矩阵元：}在多个格点间距和多个核子动量下，计算等时非局域关联函数的裸矩阵元$\tilde{h}^B(z, P^z, 1/a)$。
    \item \textbf{非微扰重整化：}应用混合重整化方案，消除线性发散和对数发散，得到$\tilde{h}^R(z, P^z)$。
    \item \textbf{大$\lambda$外推：}使用$\tilde{h}^R(\lambda) = l_1 \lambda^{-a_1} e^{-\lambda/\lambda_0}$补全大Ioffe-时间区域的矩阵元。
    \item \textbf{Fourier变换：}通过Fourier变换得到动量空间中的准PDF：
    $\tilde{x}\tilde{g}(x, P^z) = \frac{1}{\mathcal{N}}\int \frac{dz}{2\pi P^z} e^{-ixzP^z} \tilde{h}^R(z, P^z)$。
    \item \textbf{微扰匹配：}通过NLO匹配核将准PDF转换为$\overline{\text{MS}}$方案下的光锥PDF。
    \item \textbf{连续极限和无穷大动量外推：}利用式(\ref{eq:joint_extrap})消除剩余的格点 artefacts 和高阶twist污染。
\end{enumerate}

\section{光前图像与Euclidean图像的对应关系}

\subsection{时空关联长度的boost行为}

LaMET与光前量子化之间存在着深刻的时空对应关系\cite{Ji2014LaMET}。考虑一个空间关联长度$\ell$沿$z$轴：

\begin{itemize}
    \item 在LaMET图像中（高动量强子、Euclidean空间关联）：关联长度$\ell \sim 1/(xP^z)$随boost因子$\gamma$被\textbf{压缩}（Lorentz收缩）。
    \item 在光前图像中（静止强子、光锥时间关联）：同样的关联在$\xi^-$方向被boost，长度\textbf{增大}了$\gamma$倍。
\end{itemize}

这一关系如图\ref{fig:spacetime}所示（参见原文\cite{Ji2014LaMET}图1）。这种互补的视角解释了为什么：

\begin{itemize}
    \item 对于\textbf{价部分子}（$x \sim 0.1\text{--}1$），关联长度$\sim 1/\Lambda_{\text{QCD}} \sim 1$~fm，两种图像中的计算都是自然的。
    \item 对于\textbf{小$x$部分子}（$x \lesssim 10^{-2}$），关联长度$\sim P^z/(M\Lambda_{\text{QCD}})$在光前图像中变得非常长，而在LaMET中需要很大的$P^z$才能达到较小的$x$：$x \sim \Lambda_{\text{QCD}}/P^z$。
\end{itemize}

\begin{figure}[H]
\centering
\begin{tikzpicture}
    % Left panel: LaMET view
    \draw[->] (-0.5,0) -- (4,0) node[right] {$z$（空间）};
    \draw[->] (0,-0.5) -- (0,3) node[above] {$t$（虚时间）};
    \draw[thick, blue, <->] (0,0.5) -- (2,0.5);
    \node[blue, font=\small] at (1,0.8) {$\ell \sim 1/(xP^z)$};
    \node[font=\small] at (2,-1) {Euclidean关联（LaMET）};

    % Right panel: Light-front view
    \draw[->] (6,-0.5) -- (10.5,-0.5) node[right] {$\xi^-$};
    \draw[->] (7,0) -- (7,3) node[above] {$\xi^+$（光前时间）};
    \draw[thick, red, <->] (7,0.5) -- (9.5,0.5);
    \node[red, font=\small] at (8.25,0.8) {$\gamma\ell \sim 1/\Lambda_{\text{QCD}}$};
    \node[font=\small] at (8.5,-0.3) {光前关联};
\end{tikzpicture}
\caption{LaMET的空间关联与光前量子化的光锥时间关联之间的对应关系}
\label{fig:spacetime}
\end{figure}

\subsection{Minkowski空间与Euclidean空间的张量转换}

在实际计算中，正确处理Minkowski量与Euclidean量之间的转换关系至关重要\cite{NoteGluonPDFs}。以规范场张量为例：

\begin{equation}
F^{(M)}_{\mu\nu} =
\begin{pmatrix}
0   & F_{tx} & F_{ty} & F_{tz} \\
F_{xt} & 0      & F_{xy} & F_{xz} \\
F_{yt} & F_{yx} & 0       & F_{yz} \\
F_{zt} & F_{zx} & F_{zy}  & 0
\end{pmatrix}_{\!\!M}, \quad
F^{\mu\nu,(M)} =
\begin{pmatrix}
0      & -F_{tx} & -F_{ty} & -F_{tz} \\
-F_{xt} & 0       & F_{xy}  & F_{xz} \\
-F_{yt} & F_{yx}  & 0       & F_{yz} \\
-F_{zt} & F_{zx}  & F_{zy}  & 0
\end{pmatrix}_{\!\!M}。
\label{eq:F_Minkowski_full}
\end{equation}

Euclidean指标与Minkowski指标的对应关系为$\gamma_4 = \gamma^0_M$，$\gamma_i = i\gamma^i_M$（$i = 1,2,3$）。这些转换关系直接影响胶子流的构造，如式(\ref{eq:gluon_current})所示。

\subsection{LaMET与高能散射实验的类比}

季向东指出\cite{Ji2014LaMET}，LaMET与高能散射实验中提取PDF的方式有着深刻的结构类比：

\begin{center}
\begin{tabular}{lcc}
\toprule
& 高能散射实验 & LaMET/格点QCD \\
\midrule
出发点        & 依赖于大动量$Q$的散射截面 & 依赖于大强子动量$P^z$的准观测量 \\
标度分离      & 因子化定理                  & LaMET因子化 \\
联系对象      & 截面$\leftrightarrow$PDF    & 准观测量$\leftrightarrow$光锥量 \\
匹配          & 硬散射截面（微扰可计算）    & 匹配核$Z$（微扰可计算） \\
非微扰输入    & 实验数据                    & 格点``数据'' \\
\bottomrule
\end{tabular}
\end{center}

正如同一个PDF可以从不同的硬散射过程中提取，同一个光锥物理量也可以从不同的格点算符中提取。所有产生相同光锥物理的算符构成了一个\textbf{普适性类}（universality class）\cite{Ji2014LaMET}。

\section{总结与展望}

从Feynman的无穷大动量系到光锥关联函数的现代形式，再到LaMET框架的建立，我们见证了理论物理学家如何巧妙地绕过``Euclidean格点无法直接计算Minkowski光锥量''这一基本障碍。

\textbf{核心思想}可以凝练为一句话\cite{Ji2013Euclidean}：

\begin{quote}
``The light-cone correlations of quarks and gluons can be calculated by boosting the matrix elements of spatial correlations to a large momentum.''
\end{quote}

\textbf{关键里程碑}包括：

\begin{enumerate}
    \item \textbf{2013年：}季向东提出LaMET的基本思想——用大动量boost连接Euclidean空间关联与光锥关联\cite{Ji2013Euclidean}。
    \item \textbf{2014年：}LaMET被系统表述为大动量有效场论，建立了与HQET的类比\cite{Ji2014LaMET}。
    \item \textbf{2018年：}Izubuchi等人给出了准PDF与光锥PDF之间因子化定理的严格证明，并证明了准PDF与赝PDF方法的等价性\cite{Izubuchi2018}。
    \item \textbf{2020--2021年：}混合重整化方案\cite{Ji2021hybrid}和自重整化方案的提出，解决了格点上准光前关联函数的非微扰重整化问题。
    \item \textbf{2022--2023年：}Yao等人\cite{Yao2023}给出了味单态和非单态组合在非前行运动学下的完整匹配核统一框架，覆盖了GPDs、PDFs和DAs。
    \item \textbf{2024--2025年：}CLQCD/LPC合作组\cite{Chen2025gluon}在多格点间距系综上完成了核子非极化胶子PDF的连续极限计算，HadStruc合作组\cite{Egerer2022}完成了极化胶子Ioffe-时间分布的首次数值计算。
\end{enumerate}

\textbf{未来方向}包括：
\begin{itemize}
    \item 推进到更小的$x$区域（需要更大的$P^z$和更长的格点$z$方向）；
    \item 在物理$\pi$介子质量和更细格点间距下进行连续极限外推；
    \item 发展NNLO匹配核以提高微扰匹配的精度；
    \item 将重整化群重求和（RGR）和领头重正子重求和（LRR）纳入匹配程序\cite{Zhang2023RGR,Su2023LRR}；
    \item 将LaMET框架系统地推广到TMDs、GPDs和高阶twist分布。
\end{itemize}

LaMET已经从根本上改变了格点QCD中部分子物理的计算范式。随着未来EIC实验的开展和百亿亿次级（exascale）计算资源的部署，这一框架有望为实现从第一性原理精确确定核子三维结构的宏伟目标提供决定性的理论支撑。

\begin{thebibliography}{99}

% === 核心理论文献 ===

\bibitem{Ji2013Euclidean}
X.~Ji,
``Parton physics on a Euclidean lattice,''
Phys.\ Rev.\ Lett.\ \textbf{110}, 262002 (2013), arXiv:1305.1539.
\quad\textbf{[来源:} \texttt{docs/Parton Physics on Euclidean Lattice.pdf}\textbf{]}

\bibitem{Ji2014LaMET}
X.~Ji,
``Parton physics from large-momentum effective field theory,''
Sci.\ China Phys.\ Mech.\ Astron.\ \textbf{57}, 1407 (2014), arXiv:1404.6680.
\quad\textbf{[来源:} \texttt{docs/Parton Physics from Large-Momentum Effective Field Theory.pdf}\textbf{]}

\bibitem{Izubuchi2018}
T.~Izubuchi, X.~Ji, L.~Jin, I.~W.~Stewart, and Y.~Zhao,
``Factorization theorem relating Euclidean and light-cone parton distributions,''
Phys.\ Rev.\ D \textbf{98}, 056004 (2018), arXiv:1801.03917.
\quad\textbf{[来源:} \texttt{docs/Factorization Theorem Relating Euclidean and Light-Cone Parton Distributions.pdf}\textbf{]}

\bibitem{Ji2017More}
X.~Ji, J.-H.~Zhang, and Y.~Zhao,
``More on large-momentum effective theory approach to parton physics,''
Nucl.\ Phys.\ B \textbf{924}, 366 (2017), arXiv:1706.07416.
\quad\textbf{[来源:} \texttt{docs/More On Large-Momentum Effective Theory Approach to Parton Physics.pdf}\textbf{]}

\bibitem{Yao2023}
F.~Yao, Y.~Ji, and J.-H.~Zhang,
``Connecting Euclidean to light-cone correlations: From flavor nonsinglet in forward kinematics to flavor singlet in non-forward kinematics,''
JHEP \textbf{11}, 021 (2023), arXiv:2212.14415.
\quad\textbf{[来源:} \texttt{docs/Connecting Euclidean to light-cone correlations From flavor nonsinglet in forward kinematics to flavor singlet in non-forward kinematics.pdf}\textbf{]}

% === 方法论文献 ===

\bibitem{Ji2021hybrid}
X.~Ji, Y.~Liu, A.~Sch\"afer, W.~Wang, Y.-B.~Yang, J.-H.~Zhang, and Y.~Zhao,
``A hybrid renormalization scheme for quasi light-front correlations in large-momentum effective theory,''
Nucl.\ Phys.\ B \textbf{964}, 115311 (2021), arXiv:2008.03886.
\quad\textbf{[来源:} \texttt{docs/A Hybrid Renormalization Scheme for Quasi Light-Front Correlations in Large-Momentum Effective Theory.pdf}\textbf{]}

\bibitem{Zhang2019}
J.-H.~Zhang, X.~Ji, A.~Sch\"afer, W.~Wang, and S.~Zhao,
``Accessing gluon parton distributions in large momentum effective theory,''
Phys.\ Rev.\ Lett.\ \textbf{122}, 142001 (2019), arXiv:1808.10824.
\quad\textbf{[来源:} \texttt{docs/Accessing gluon parton distributions in large momentum effective theory.pdf}\textbf{]}

\bibitem{Balitsky2020}
I.~Balitsky, W.~Morris, and A.~Radyushkin,
``Short-distance structure of unpolarized gluon pseudodistributions,''
Phys.\ Lett.\ B \textbf{808}, 135621 (2020), arXiv:1910.13963.
\quad\textbf{[来源:} \texttt{docs/Short-Distance Structure of Unpolarized Gluon Pseudodistributions.pdf}\textbf{]}

\bibitem{Chen2017Renorm}
J.-W.~Chen, X.~Ji, and J.-H.~Zhang,
``A complete non-perturbative renormalization prescription for quasi-PDFs,''
Nucl.\ Phys.\ B \textbf{915}, 1 (2017), arXiv:1609.08102.
\quad\textbf{[来源:} \texttt{docs/A complete non-perturbative renormalization prescription for quasi-PDFs.pdf}\textbf{]}

\bibitem{Alexandrou2017}
C.~Alexandrou, K.~Cichy, M.~Constantinou, K.~Hadjiyiannakou, K.~Jansen, H.~Panagopoulos, and F.~Steffens,
``A complete non-perturbative renormalization prescription for quasi-PDFs,''
Nucl.\ Phys.\ B \textbf{923}, 394 (2017), arXiv:1706.00265.

\bibitem{Li2019Renorm}
Z.-Y.~Li, Y.-Q.~Ma, and J.-W.~Qiu,
``Self-renormalization of quasi-light-front correlators on the lattice,''
Phys.\ Rev.\ Lett.\ \textbf{122}, 062002 (2019), arXiv:1809.01836.
\quad\textbf{[来源:} \texttt{docs/Self-Renormalization of Quasi-Light-Front Correlators on the Lattice.pdf}\textbf{]}

% === 数值应用文献 ===

\bibitem{Chen2025gluon}
C.~Chen \textit{et al.} (CLQCD, Lattice Parton),
``Unpolarized gluon PDF of the nucleon from lattice QCD in the continuum limit,''
(2025), arXiv:2510.26425.
\quad\textbf{[来源:} \texttt{docs/Unpolarized gluon PDF of the nucleon from lattice QCD in the continuum limit.pdf}\textbf{]}

\bibitem{Chen2025first}
C.~Chen \textit{et al.} (CLQCD, Lattice Parton),
``First self-renormalized gluon PDF of nucleon from large-momentum effective theory in the continuum limit,''
Phys.\ Rev.\ D \textbf{111}, 074506 (2025), arXiv:2408.12819.
\quad\textbf{[来源:} \texttt{docs/First Self-Renormalized Gluon PDF of Nucleon from Large-Momentum Effective Theory in the Continuum Limit.pdf}\textbf{]}

\bibitem{Egerer2022}
C.~Egerer \textit{et al.} (HadStruc Collaboration),
``Towards the determination of the gluon helicity distribution in the nucleon from lattice quantum chromodynamics,''
Phys.\ Rev.\ D \textbf{106}, 094511 (2022), arXiv:2207.08733.
\quad\textbf{[来源:} \texttt{docs/Towards the determination of the gluon helicity distribution in the nucleon from lattice quantum chromodynamics.pdf}\textbf{]}

\bibitem{LP3helicity}
H.-W.~Lin, J.-W.~Chen, X.~Ji, L.~Jin, R.~Li, Y.-S.~Liu, Y.-B.~Yang, J.-H.~Zhang, and Y.~Zhao (LP$^3$ Collaboration),
``Proton isovector helicity distribution on the lattice at physical pion mass,''
Phys.\ Rev.\ Lett.\ \textbf{121}, 242003 (2018), arXiv:1807.07431.
\quad\textbf{[来源:} \texttt{docs/Proton Isovector Helicity Distribution on the Lattice at Physical Pion Mass.pdf}\textbf{]}

\bibitem{Good2025a}
W.~Good, F.~Yao, and H.-W.~Lin,
``First nucleon gluon PDF from large momentum effective theory,''
(2025), arXiv:2505.13321.
\quad\textbf{[来源:} \texttt{docs/First Nucleon Gluon PDF from Large Momentum Effective Theory.pdf}\textbf{]}

\bibitem{Lin2025gluon}
W.~Good, K.~Hasan, and H.-W.~Lin,
``Gluon parton distribution of the nucleon from $2+1+1$-flavor lattice QCD in the physical-continuum limit,''
J.\ Phys.\ G \textbf{52}, 035105 (2025), arXiv:2409.02750.
\quad\textbf{[来源:} \texttt{docs/Gluon Parton Distribution of the Nucleon from 2+1+1-Flavor Lattice QCD in the Physical-Continuum Limit.pdf}\textbf{]}

\bibitem{NoteGluonPDFs}
``Note of gluon PDFs,''
内部笔记。（含Minkowski/Euclidean张量转换关系的详细推导）
\quad\textbf{[来源:} \texttt{docs/Note of gluon PDFs.pdf}\textbf{]}

% === 改进方法 ===

\bibitem{Zhang2023RGR}
R.~Zhang, J.~Holligan, X.~Ji, and Y.~Su,
``Resumming quark's longitudinal momentum logarithms in LaMET expansion of lattice PDFs,''
Phys.\ Lett.\ B \textbf{844}, 138081 (2023), arXiv:2305.05212.
\quad\textbf{[来源:} \texttt{docs/Resumming Quark's Longitudinal Momentum Logarithms in LaMET Expansion of Lattice PDFs.pdf}\textbf{]}

\bibitem{Su2023LRR}
Y.~Su, J.~Holligan, X.~Ji, F.~Yao, J.-H.~Zhang, and R.~Zhang,
``Power corrections and renormalons in parton quasi-distributions,''
Nucl.\ Phys.\ B \textbf{991}, 116201 (2023), arXiv:2209.01236.
\quad\textbf{[来源:} \texttt{docs/Power corrections and renormalons in parton quasi-distributions.pdf}\textbf{]}

% === 其他相关文献 ===

\bibitem{CLQCD2024}
Z.-C.~Hu \textit{et al.} (CLQCD),
``Charmed meson masses and decay constants in the continuum from the tadpole improved clover ensembles,''
Phys.\ Rev.\ D \textbf{109}, 054507 (2024), arXiv:2310.00814.
\quad\textbf{[来源:} \texttt{docs/Charmed meson masses and decay constants in the continuum from the tadpole improved clover ensembles.pdf}\textbf{]}

\bibitem{CLQCD2025}
H.-Y.~Du \textit{et al.} (CLQCD),
``Quark masses and low energy constants in the continuum from the tadpole improved clover ensembles,''
Phys.\ Rev.\ D \textbf{111}, 054504 (2025), arXiv:2408.03548.
\quad\textbf{[来源:} \texttt{docs/Quark masses and low energy constants in the continuum from the tadpole improved clover ensembles.pdf}\textbf{]}

\bibitem{Fan2023}
Z.~Fan, W.~Good, and H.-W.~Lin,
``Physical-continuum limit of the nucleon gluon parton distribution from lattice QCD,''
Phys.\ Rev.\ D \textbf{108}, 014508 (2023), arXiv:2210.09985.
\quad\textbf{[来源:} \texttt{docs/Physical-Continuum Limit of the Nucleon Gluon Parton Distribution from Lattice QCD.pdf}\textbf{]}

\bibitem{Zhang2025kinematic}
R.~Zhang, A.~V.~Grebe, D.~C.~Hackett, M.~L.~Wagman, and Y.~Zhao,
``Kinematically-enhanced interpolating operators for boosted hadrons,''
(2025), arXiv:2501.00729.
\quad\textbf{[来源:} \texttt{docs/Kinematically-enhanced interpolating operators for boosted hadrons.pdf}\textbf{]}

\bibitem{Egerer2021a}
C.~Egerer, R.~G.~Edwards, K.~Orginos, and D.~G.~Richards,
``Distillation at high-momentum,''
Phys.\ Rev.\ D \textbf{103}, 034502 (2021), arXiv:2009.10691.
\quad\textbf{[来源:} \texttt{docs/Distillation at High-Momentum.pdf}\textbf{]}

\bibitem{Peardon2009}
M.~Peardon \textit{et al.} (Hadron Spectrum),
``A novel quark-field creation operator construction for hadronic physics in lattice QCD,''
Phys.\ Rev.\ D \textbf{80}, 054506 (2009), arXiv:0905.2160.
\quad\textbf{[来源:} \texttt{docs/A novel quark-field creation operator construction for hadronic physics in lattice QCD.pdf}\textbf{]}

\end{thebibliography}

\end{document}
